% Generated from validated Article Draft Result. Do not hand-edit; revise the structured draft instead.
\section{振る舞いを教える――instruction、feedback、verifierが分かれ始めた}
\label{pkg:instruction-human-feedback}
\noindent\textbf{FLAN、T0、human-feedback summarization、verifierは、すべて『モデルを賢くする追加学習』として一括できるものではない。instructionへの適応、multitask prompted training、人間の評価を用いたsteering、候補を評価するverifierという別々の操作点が現れ、pretraining後の振る舞いを設計する層が2021年に輪郭を持ち始めた。}\autocite{src-34f845c1e710,src-6339600396eb,src-750bc3eeb1bc,src-7f768525daf5,src-8b2f7c40f4a4,src-90a37f62aeab}

2021年のsteeringを理解するには、promptを入力時に工夫することと、instructionを使って学習段階でモデルを変えることを分ける必要がある。FLANが代表するinstruction finetuningは、pretrained modelを指示へ合わせる後段の学習を独立した問題として見せた。この境界は、前章のparameter-efficient adaptationと同じく、pretraining済み基盤を固定的な完成品ではなく、その後に振る舞いを設計できるsubstrateとして扱う流れに位置付く。\autocite{src-8b2f7c40f4a4}


T0は、multitask prompted trainingという別の設計軸を示す。複数タスクをprompt形式で扱う学習と、単一の入力に対して実行時promptだけを調整することは同じではない。2021年には『prompt』という語が入力interfaceだけでなくtraining designにも入り込み、どの段階で、何を教師信号として、どの範囲の振る舞いを整えるかという問いが分化していった。\autocite{src-6339600396eb}


Recursive Book Summarization with Human Feedbackは、長い対象を扱うsummarizationで、人間から得た評価をsteeringへ組み込む系統を記録している。ここで年次的に重要なのは、正解ラベルだけを与える学習とは別に、人間が望む出力の評価をモデル改善へ接続する設計が明示的な研究対象になっていた点である。後年の手法を逆投影せず、2021年時点ではlong-horizon summarizationにおけるhuman-feedback steeringの一次資料として位置付ける。\autocite{src-7f768525daf5,src-90a37f62aeab}


Training Verifiers / GSM8Kは、生成そのものだけでなく、候補を評価するverifierを組み合わせる方向を示した。本稿ではこれを2021年に完成した『reasoning system』とは呼ばない。むしろ、生成器の一回の出力をそのまま最終結果とせず、評価側のモデルや信号を追加できるという構造が、後のreasoning-oriented systemsへつながる橋として観測できる、と限定して読む。\autocite{src-34f845c1e710,src-750bc3eeb1bc}


四つの系統を並べると、pretraining後の操作点が一つではなかったことが分かる。instruction finetuningは指示への適応、multitask prompted trainingはタスク横断の学習構成、human feedbackは人間評価との接続、verifierは生成結果を選別・評価する補助系である。2021年の転換は、巨大モデルの能力を測るだけでなく、その能力をどう引き出し、整え、評価するかが独立した設計空間になったことにある。\autocite{src-34f845c1e710,src-6339600396eb,src-750bc3eeb1bc,src-7f768525daf5,src-8b2f7c40f4a4,src-90a37f62aeab}


\begin{claimboundary}
各研究が報告する性能や有効性は原著者の評価であり、ここでは独立再現済みの事実として扱わない。また、2022年以降に一般化したinstruction-followingやreasoningの完成形を2021年の成果へ遡及的に読み込まない。\autocite{src-34f845c1e710,src-6339600396eb,src-750bc3eeb1bc,src-7f768525daf5,src-8b2f7c40f4a4,src-90a37f62aeab}
\end{claimboundary}
